The algebra is finite-dimensional, so with any norm satisfying $\lVert xy \rVert \le c \lVert x \rVert \lVert y \rVert$ the series $\sum_k X^{k}/k!$ converges absolutely for every $X$, and products of two such series may be multiplied term by term (Cauchy product). Nothing below uses more analysis than this, except the surjectivity of the complex hyperbolic cosine in (iv).

(i) Each $B^{k}$ is a product of even elements, so $e^{B}$ is even. Reversion is a linear anti-automorphism with $\tilde B = -B$ on bivectors (\claimref{std-algebra.involution-properties}), so $(B^{k})^{\sim} = \tilde B^{k} = (-B)^{k}$ and, reversion being linear on a finite-dimensional space, $(e^{B})^{\sim} = e^{-B}$. Since $sB$ and $tB$ commute, the Cauchy product gives $e^{sB} e^{tB} = \sum_m \sum_{j + k = m} s^{j} t^{k} B^{m}/(j!\,k!) = \sum_m (s + t)^{m} B^{m}/m! = e^{(s+t)B}$ by the binomial theorem; with $s = 1$, $t = -1$ this is $e^{B}e^{-B} = 1$. Hence $e^{B}(e^{B})^{\sim} = 1$. None of this uses the field or the form.

(ii) Let $L_B(y) = By$ and $R_B(y) = yB$. These linear maps of the algebra commute ($B(yB) = (By)B$), so $(L_B - R_B)^{m} = \sum_{j} \binom{m}{j} L_B^{j}(-R_B)^{m-j}$, that is $\operatorname{ad}_B^{m}(y) = \sum_j \binom{m}{j} B^{j} y (-B)^{m-j}$. Multiplying the series for $e^{B}$ and $e^{-B}$ with $y$ between them and collecting terms of total degree $m$,
\[
e^{B} y e^{-B} = \sum_{m} \sum_{j + k = m} \frac{B^{j} y (-B)^{k}}{j!\,k!} = \sum_m \frac{\operatorname{ad}_B^{m}(y)}{m!} .
\]
By \claimref{std-algebra.bivector-lie-algebra}, $\operatorname{ad}_B(x) = 2 B \times x = 2\langle Bx \rangle_1 \in V$ for $x \in V$, so every $\operatorname{ad}_B^{m}(x)$ is a vector, and so is the convergent sum, $V$ being a closed subspace. The map $x \mapsto e^{B} x e^{-B}$ is linear and $(e^{B} x e^{-B})^{2} = e^{B} x^{2} e^{-B} = Q(x)$, so it preserves $Q$; it is invertible with inverse given by $-B$. Over $\C$, or for degenerate forms, \claimref{std-algebra.bivector-lie-algebra} (stated for real spaces) is not needed: the defining relation $v^{2} = Q(v)$, polarised, gives $xu + ux = 2B(u, x)$ in any Clifford algebra over $\R$ or $\C$, so for vectors $u, v, x$, $(uv)x - x(uv) = u(vx + xv) - (ux + xu)v = 2B(v, x)\,u - 2B(u, x)\,v \in V$; every bivector is a sum of terms $uv - B(u, v)$ (over $\C$ by the blade basis of \claimref{std-algebra.complex-clifford-algebra}(i)), so $\operatorname{ad}_B(V) \subseteq V$ for every $B$, over either field and for every form.

(iii) Let $Q$ be real and nondegenerate, and $M(x) = e^{B} x e^{-B}$, an isometry by (ii). By \claimref{std-algebra.cartan-dieudonne}, $M = s_{u_1} \cdots s_{u_k}$ with non-null $u_i$. Put $g = u_1 \cdots u_k$; by \claimref{primer.twisted-adjoint-action}, $\Ad_g = \Ad_{u_1} \cdots \Ad_{u_k} = s_{u_1}\cdots s_{u_k} = M$, where $\Ad_g(x) = \hat g x g^{-1}$. Since $e^{B}$ is even, $\Ad_{e^{B}} = M$ too, so $h = g^{-1}e^{B}$, which is homogeneous of the parity of $k$ and invertible, satisfies $\hat h x h^{-1} = x$ for all $x \in V$. Take an orthogonal basis $e_i$ with $Q(e_i) = \pm 1$. If $k$ is odd, $h$ is odd and $h x = -x h$ for every vector; conjugation by $e_j$ multiplies a blade $e_I$ by $+1$ or $-1$ according as $j$ is outside or inside $I$ when $\lvert I \rvert$ is even, and by $-1$ or $+1$ when $\lvert I \rvert$ is odd; so each blade $e_I$ occurring in an odd $h$ that anticommutes with every $e_j$ would need $j \notin I$ for all $j$, i.e. $I = \emptyset$, which is not odd; hence $h = 0$, impossible. So $k$ is even, $h$ is even and commutes with every vector, hence is central; the centre is $\R$ for $n$ even and $\R \oplus \R\omega$ with $\omega$ odd for $n$ odd (\claimref{primer.centre}), so $h = \lambda \in \R \setminus \{0\}$. Thus $e^{B} = \lambda u_1 \cdots u_k$ with $k$ even, and by (i)
\[
1 = e^{B} (e^{B})^{\sim} = \lambda^{2}\, u_1 \cdots u_k u_k \cdots u_1 = \lambda^{2} \prod_i Q(u_i),
\]
so $\prod_i Q(u_i) > 0$. Replace $u_i$ by $u_i / \sqrt{\lvert Q(u_i) \rvert}$ and absorb $\lambda \prod_i \sqrt{\lvert Q(u_i)\rvert} = \pm 1$ into $u_1$: then $e^{B} = u_1 \cdots u_k$ with $Q(u_i) = \pm 1$ and $\prod Q(u_i) = +1$. So $e^{B} \in \Spin^{+}(V, Q)$ (\claimref{obj.std-algebra.pin-spin}), it is an even versor with reverse-norm $1$, a rotor, and its action $M$ is an even product of reflections with $\prod Q(u_i) > 0$, which is $\mathrm{SO}^{+}(V, Q)$ by definition (determinant $1$ by \claimref{std-algebra.spin-double-cover}). Nondegeneracy is needed: in $\Cl(0,0,2)$ every vector $v$ has $Q(v) = 0$, so no vector is invertible (\claimref{std-algebra.inverse-of-vector}) and the only versor is the empty product $1$; but $B = f_1 f_2$ has $B^{2} = 0$, so $e^{B} = 1 + f_1 f_2$, which has $e^{B}(e^{B})^{\sim} = 1$ and fixes every vector, and is not a versor.

(iv) In $\Cl(1,3)$, $n = 4$ is even, so $I$ anticommutes with vectors and commutes with even elements (\claimref{primer.centre}), and $I^{2} = -1$ (\claimref{primer.sign}). So $\C_I$ is a commutative subalgebra with $(a + bI)(a - bI) = a^{2} + b^{2}$, a field isomorphic to $\C$, central in $\Cl^{+}$. For a bivector $B$, $B^{2} = \langle B^{2}\rangle_0 + \langle B^{2}\rangle_2 + \langle B^{2}\rangle_4$ and $\langle B^{2}\rangle_2 = B \times B = 0$, so $z := B^{2} \in \C_I$. Then $B^{2k} = z^{k}$ and $B^{2k+1} = z^{k} B$, so
\[
e^{B} = C(z) + S(z)\, B, \qquad C(z) = \sum_k \frac{z^{k}}{(2k)!}, \quad S(z) = \sum_k \frac{z^{k}}{(2k+1)!},
\]
entire functions on $\C_I$ with $C(w^{2}) = \cosh w$, $S(w^{2}) = \sinh w / w$ ($S(0) = 1$) and $C(z)^{2} - z S(z)^{2} = 1$; for $z \in \R$ these are the closed forms of \claimref{std-algebra.bivector-exponential}. An even $R$ has grades $0, 2, 4$, so $R = s + F$ with $s \in \C_I$ and $F$ a bivector; $\tilde R = s - F$ ($\tilde I = I$), and $R\tilde R = s^{2} - F^{2}$ since $s$ commutes with $F$. So $R\tilde R = 1$ reads $s^{2} - F^{2} = 1$ in $\C_I$.

Suppose $s \ne -1$. If $s = 1$ then $F^{2} = 0$ and $R = 1 + F = e^{F}$. Otherwise choose $w \in \C_I$ with $\cosh w = s$ (the complex hyperbolic cosine is onto); then $\sinh^{2} w = s^{2} - 1 \ne 0$, so $S(w^{2}) \ne 0$, and $B := F / S(w^{2})$ has $B^{2} = F^{2}/S(w^{2})^{2} = (s^{2} - 1) w^{2}/\sinh^{2} w = w^{2}$; hence $e^{B} = \cosh w + S(w^{2}) B = s + F = R$. If $s = -1$ and $F = 0$, $R = -1 = e^{\pi \gamma_1\gamma_2}$, since $(\gamma_1\gamma_2)^{2} = -1$. If $s = -1$ and $F = N \ne 0$, then $F^{2} = s^{2} - 1 = 0$; suppose $e^{B} = R$. Comparing the $\C_I$ and bivector parts, $C(z) = -1$ and $S(z) B = N \ne 0$. Write $z = w^{2}$: $\cosh w = -1$ forces $w = (2m + 1)\pi I$ for an integer $m$, where $\sinh w = 0$ and $w \ne 0$, so $S(z) = \sinh w / w = 0$, a contradiction. So $-(1 + N)$ is not an exponential, while $-R = 1 + N = e^{N}$; every such $R$ is therefore $\pm e^{B}$. Finally $N_0 = (\gamma_0 + \gamma_1)\gamma_2$ has $N_0^{2} = 0$ (the null vector $\gamma_0 + \gamma_1$ is orthogonal to $\gamma_2$), and $-(1 + N_0) = (\gamma_1\gamma_3\gamma_1\gamma_3)(-\gamma_2)(\gamma_2 + \gamma_0 + \gamma_1)$ is a product of six vectors of square $-1$, so it is a rotor that is not an exponential.
